\begin{abstract}
Key-value (KV) caching is essential for efficient autoregressive Transformer inference, but its memory footprint grows linearly with context length, batch size, number of layers, and number of attention heads. Existing KV cache compression methods typically decide what to retain from attention statistics, recency, layer structure, or quantization error. These signals are useful, but they do not explicitly account for when compression-induced changes are likely to be unsafe. This draft proposes \uckv{}, an uncertainty-calibrated framework for adaptive KV cache compression. \uckv{} treats cache compression as a selective memory decision: a calibrated risk model predicts whether a candidate cache budget will perturb the next-token distribution, and the system chooses the smallest budget below a risk tolerance, falling back to conservative retention when needed. We provide a formal problem statement, an attention perturbation bound linking discarded attention mass to output deviation, and a calibration-based risk-control argument. We also report controlled synthetic retrieval pilots using GPT-2 and SmolLM2-135M-Instruct. In the final SmolLM2 pilot, the selected tolerance $\epsilon=0.36$ matches the answer-containment accuracy of full-cache decoding and the strongest fixed sink-plus-window baseline on two held-out confirmation seeds, while reducing average retained tokens from 286.07 to 266.10 relative to full cache. These results are preliminary: they support uncertainty-calibrated budget selection as a promising operating point, but they do not yet establish wall-clock speedups or state-of-the-art long-context performance.
\end{abstract}
